% !TEX root = ../Report.tex
\chapter*{Abstract}
\addcontentsline{toc}{chapter}{Abstract}

Trip planning becomes significantly more complex when multiple people collaborate,
because itineraries evolve through parallel edits, last-minute changes, and
disagreements about priorities. Existing tools typically optimise a single
itinerary (e.g., mapping and routing applications) or support generic
collaboration (e.g., shared documents), but they do not provide robust
change-management semantics tailored to itineraries. This project presents
\textit{GiTrip}, a server-rendered web application that applies Git-style
version control concepts to trip planning: a trip is modelled as a repository,
alternative itineraries are branches, each saved version is a commit containing
a complete snapshot, and divergent branches can be merged using a three-way
merge with an explicit conflict-resolution user interface.

GiTrip integrates routing, scheduling constraints, and usability-focused
interaction. An automatic scheduler produces feasible multi-day itineraries
based on user inputs such as active hours, stay durations, ordering constraints,
desired time windows, public transit selection, and opening hours. Users can
then refine plans using two interactive surfaces: a timeline editor for
adjusting arrival/departure times and a map editor for repositioning locations.
To address the usability barrier of Git concepts for non-technical users, GiTrip
includes an ``Easy mode'' that hides advanced controls while preserving version
safety.

The completed system was evaluated through an in-person Evaluation Day
demonstration and user acceptance testing. Objective task results show high
completion rates for auto-scheduling feasibility and merge conflict resolution,
with qualitative feedback highlighting the novelty and usefulness of
version-controlled planning. The report details the requirements, system design,
implementation, testing strategy, and evaluation outcomes, and concludes with
limitations and future work (including optional machine learning extensions
beyond the original project scope).

\vspace{0.6em}
\textbf{Keywords:} trip planning; version control; branching and merging;
scheduling constraints; routing; usability; progressive web app; conflict
resolution
